% Part: turing-machines 
% Chapter: undecidability
% Section: introduction

\documentclass[../../../include/open-logic-section]{subfiles}

\begin{document}

\olfileid[pt]{tur}{und}{int} 
\olsection{Introdução}

Pode parecer óbvio que nem toda função, nem mesmo toda função
aritmética, pode ser computável. Há simplesmente um número grande
demais delas, cujo comportamento é complicado demais. Funções
definidas a partir do decaimento de partículas radioativas, por
exemplo, ou de outro comportamento caótico ou aleatório. Suponha que
comecemos a contar intervalos de 1 segundo a partir de um dado
instante e definamos a função $f(n)$ como o número de partículas no
universo que decaem no $n$-ésimo intervalo de 1 segundo após aquele
momento inicial. Esta parece ser uma candidata a função que jamais
poderemos ter esperança de computar.

Mas uma coisa é não conseguir imaginar como se computaria tais
funções, e outra bem diferente é de fato provar que elas são
incomputáveis. Na verdade, mesmo funções que parecem irremediavelmente
complicadas podem, em um sentido abstrato, ser computáveis. Por
exemplo, suponha que o universo seja finito no tempo---algum dia, num
futuro muito distante, o universo se contrairá em um único ponto, como
preveem algumas teorias cosmológicas. Então há apenas um número finito
(mas incrivelmente grande) de segundos a partir daquele momento
inicial para os quais $f(n)$ está definida. E qualquer função que
esteja definida para apenas finitas entradas é computável: poderíamos
listar as saídas em uma grande tabela, ou codificá-la em um diagrama
de transição de estados de máquina de Turing muito grande.

Frequentemente estamos interessados em casos especiais de funções
cujos valores fornecem as respostas a perguntas de sim/não. Por
exemplo, a pergunta ``$n$ é um número primo?'' está associada à função
\[
\fn{primo}(n) = \begin{cases}
  1 & \text{se $n$ é primo}\\
  0 & \text{caso contrário.}
  \end{cases}
\]
Dizemos que uma pergunta de sim/não pode ser \emph{efetivamente
decidida} se a função associada de valores $1/0$ é efetivamente
computável.

Para provar matematicamente que há funções que não podem ser
efetivamente computadas, ou problemas que não podem ser efetivamente
decididos, é essencial fixar um modelo específico de computação e
mostrar que há funções que ele não pode computar ou problemas que ele
não pode decidir. Podemos mostrar, por exemplo, que nem toda função
pode ser computada por máquinas de Turing, e nem todo problema pode ser
decidido por máquinas de Turing. Podemos então apelar para a tese de
Church-Turing para concluir que não apenas as máquinas de Turing não
são poderosas o suficiente para computar toda função, como também
nenhum procedimento efetivo o é.

A chave para provar tais resultados negativos é o fato de que podemos
atribuir números às próprias máquinas de Turing. A maneira mais fácil
de fazer isso é enumerá-las, talvez fixando um modo específico de
escrever máquinas de Turing e seus programas e, em seguida,
listando-os de maneira sistemática. Uma vez que vejamos que isso pode
ser feito, então a existência de funções Turing-incomputáveis decorre
de simples considerações de cardinalidade: o conjunto de funções de
$\Nat$ a~$\Nat$ (na verdade, mesmo apenas de $\Nat$ a $\{0, 1\}$) é
!!{nonenumerable}, mas como podemos enumerar todas as máquinas de
Turing, o conjunto de funções Turing-computáveis é apenas
!!{denumerable}.

Também podemos definir funções e problemas \emph{específicos} que
podemos provar ser incomputáveis e indecidíveis, respectivamente. Um
desses problemas é o assim chamado \emph{Problema da Parada.} Máquinas
de Turing podem ser finitamente descritas listando-se suas instruções.
Tal descrição de uma máquina de Turing, isto é, um programa de máquina
de Turing, pode é claro ser usada como entrada para outra máquina de
Turing. Assim, podemos considerar máquinas de Turing que decidem
perguntas sobre outras máquinas de Turing. Uma pergunta
particularmente interessante é esta: ``A máquina de Turing dada
eventualmente para quando iniciada na entrada~$n$?'' Seria bom se
houvesse uma máquina de Turing que pudesse decidir essa pergunta:
pense nela como uma máquina de Turing de controle de qualidade que
garante que as máquinas de Turing não fiquem presas em laços infinitos
e coisas do tipo. O fato interessante, que Turing provou, é que não
pode haver tal máquina de Turing. Não pode haver uma única máquina de
Turing que, quando iniciada em uma entrada que consiste em uma
descrição de uma máquina de Turing $M$ e algum número~$n$, sempre pare
com saída $1$ ou $0$ conforme a máquina $M$ teria parado ou não quando
iniciada na entrada $n$.

Uma vez que tenhamos exemplos de problemas indecidíveis específicos,
podemos usá-los para mostrar que outros problemas também são
indecidíveis. Por exemplo, um célebre problema indecidível é a
pergunta ``A !!{formula}~$!A$ de primeira ordem é válida?''. Não há
máquina de Turing que, dada como entrada uma !!{formula}~$!A$ de
primeira ordem, tenha garantia de parar com saída $1$ ou $0$ conforme
$!A$ seja válida ou não. Historicamente, a questão de encontrar um
procedimento para resolver efetivamente esse problema foi chamada
simplesmente de ``o'' problema da decisão; e assim dizemos que o
problema da decisão é insolúvel. Turing e Church provaram esse
resultado independentemente mais ou menos na mesma época, de modo que
ele também é chamado de Teorema de Church-Turing.

\end{document}
